\cotefile{136}
\scanpage{17}\mspage{15}
\deleted{\unsure{5)|g)}}\unsure{6)|h)} On peut associer à un espace topologique $X$ \deleted{\unsure{bien}} ou un ens $\frac12$simplicial $X_*$ \deleted{\unsure{des}}divers invariants alg. ``linéaires'' \uline{plus fins} à priori que le $C^{\cdot\cdot}_{\mathrm{DRpd}}$ (ou \unsure{$=$ que|mieux par} $C_{\cdot\cdot}^{\mathrm{DRpd}}$, \unsure{j'en dis autant|j'en dirais autant|en dualisant}\ldots).

[P. ex. on peut observer que $C^{\cdot\cdot}_{\mathrm{DRpd}\,*}(k)\simeq_d\Gamma^{\cdot}\Phi_{*k}\otimes_k\bigwedge^{\cdot}\Psi_{*k}$ se déduit de $D^{\cdot\cdot}_*(k)\unsure{=|/|\overset{\mathrm{déf}}{=}}\Gamma^{\cdot}\Phi_{*k}\otimes\bigwedge^{\cdot}\Phi_{*k}$ (qui est une Alg. diff \inserted{\unsure{homogène|$k$-augm.}} à p. div. qui est une \uline{résolution de} $k$) et de $T\in\Gamma(D_*^{10}{}_{k})$ comme Coker de la multiplication par $dT$ (i.e. on divise par l'idéal engendré par $dT$), ou encore en bidegré donné \deleted{\unsure{$p,q+1$}} $(p,q)$
\[
\deleted{\unsure{D}}C^{p,q}_{*k}\simeq\operatorname{Ker}\bigl(D_*^{p,q\deleted{\unsure{+1}}}\xrightarrow[\text{produit par }dT]{}D_*^{p,q+2}\bigr)
\]
d'où
\[
C^{p,q}_{\mathrm{DRpd}}(X_*,k)\simeq\operatorname{Ker}\bigl(D^{p,q+1}(X_*,k)\unsure{-|\longrightarrow}D^{p,q+2}(X_*,k)\bigr),
\]
et on peut considérer les structures \inserted{de base} sur $C^{\cdot\cdot}_{\mathrm{DRpd}}(X_*,k)$ comme \unsure{déduite|déduites} de certaines structures (à expliciter\ldots) sur $D^{\cdot\cdot}(X_*,k)$.

On peut aussi définir ``\uuline{La} structure multiplicative à p.d \inserted{cohomologique} du type d'homotopie $X_*$'', en \inserted{un sens convenable}, qui permet de reconstituer \deleted{\unsure{les}} \inserted{aussi bien} $C^{\cdot\cdot}_{\mathrm{DRpd}}(X_*,k)$ que $D^{\cdot\cdot}_k(X_*,k)$ (ou les complexes de De Rham-Sullivan). J'ignore si \unsure{le|la|les}
% archive sheet number: 16

\scanpage{18}\mspage{16}
complexe de De Rham à p.d. \inserted{(ou \unsure{sinon|bien|même} $D^{\cdot\cdot}(X_*,\mathbb{Z})$)} suffit déjà pour récupérer toute la structure mult. \& p.d. coh. de $X_*$ \inserted{(sous réserve de conditions de finitude bien sûr)}. Dans le cas contraire, ce serait cette dernière qui serait le candidat \inserted{alg. ``linéaire''} naturel pour exprimer le type d'homotopie $X_*$ (du moins si $X_*$ \inserted{connexe et} simplement connexe, et en passant à une cat. dérivée convenable bien sûr).
% archive sheet number: 17

\scanpage{20}
\hfill\fbox{Lemmes sur coeff. binomiaux.}

\uline{Lemme} 1 Soient $p$ nb premier, $i,j\in\mathbb{N}$, avec \unsure{des développements|du développement} $p$-\unsure{adiques|adique}
\[
i=\sum_{0\leq\varepsilon\leq{e}}u_\varepsilon p^\varepsilon
\qquad j=\sum_{0\leq\varepsilon\leq e}v_\varepsilon p^\varepsilon
\qquad 0\leq u_\varepsilon,v_\varepsilon\leq p-1,
\]
soit
\[
c_{ij}=\frac{(i+j)!}{i!\,j!}\in\mathbb{Z}\qquad\text{(coeff. binômial)}
\]
alors on a
\[
c_{ij}\not\equiv0\ (p)\quad\text{sss}\quad u_\varepsilon+v_\varepsilon\leq p-1\quad\forall\varepsilon\quad(0\leq\varepsilon\leq\unsure{\nu|v})
\]
et \deleted{\unsure{pour}} on a en tous cas
\[
c_{ij}\equiv\prod_{0\leq\varepsilon\leq\nu}c_{u_\varepsilon,v_\varepsilon}\qquad(p)
\]

\margin{\uline{Cor} 1}% slanted
On écrit d'abord $i=u+pr,\quad j=v+ps\qquad 0\leq u,v\leq p-1,\quad r,s\geq0$ et on montre que
\[
\left\{\begin{aligned}
&c_{ij}\not\equiv0\ (p)\quad\text{sss}\quad u+v\leq p-1\quad\text{et}\quad c_{r,s}\not\equiv0\quad(p)\\
&c_{ij}\equiv c_{u,v}\,c_{r,s}\qquad(p)
\end{aligned}\right.\tag*{$(*)$}
\]
en étudiant le développement, dans $\mathbb{F}_p[V,Y]$, de $(X+Y)^{i+j}$. \deleted{\unsure{ss}}On a
\[
i+j=p\underbrace{(r+s)}_{t}+\underbrace{u+v}_{\deleted{\unsure{\lambda}}w}
\]
si $w\leq p-1$ on écrit
\[
\begin{aligned}
(X+Y)^{\unsure{j|i}+j}&=\bigl((X+Y)^p\bigr)^{t\deleted{\unsure{r+s}}}(X+Y)^{w}=(X^p+Y^p)^{t\deleted{\unsure{r+s}}}(X+Y)^w\\
&=\sum_{\substack{0\leq\alpha\leq w\\0\leq\beta\leq t}}\binom{w}{\alpha}\binom{t}{\beta}X^{\alpha+p\beta}Y^{(w-\alpha)+p(t-\beta)},
\end{aligned}
\]
où tous les monômes \unsure{écrit|écrits} dans la somme sont à exposants distincts. On a donc comme coeff. de $X^iY^j$ celui correspondant au terme $i=\alpha+p\beta$ i.e.
\[
\alpha=u,\quad\beta=r\qquad\text{i.e.}
\]
\[
c_{i,j}\equiv\binom{w}{u}\binom{t}{r}=c_{u,v}\,c_{r,s}\qquad(p)
\]
et il est immédiat que $c_{u,v}\not\equiv0\ (p)$, car $c_{u,v}=\frac{(u+v)!}{u!\,v!}$, $0\leq u,v,u+v\leq p-1$. Donc $c_{ij}\neq0\ (p)\iff c_{\unsure{v|r},s}\neq0\ (p)$.

Il reste à voir que si $w\geq p$, \deleted{soit $w=p+w'$} \inserted{alors} $c_{i,j}\equiv0\ (p)$. Soit $w^{\unsure{1|\prime|{}}}=p+w'\quad 0\leq w'\leq p-2$, donc $i+j=(r+s+1)p+w'=\underbrace{(t+1)}_{t'}p+w'$,
\[
(X+Y)^{i+j}=(X^p+Y^p)^{t'\deleted{\unsure{r}}}(X+Y)^{w'}=\sum_{\substack{0\leq\alpha\leq w'\\0\leq\beta\leq t'}}\binom{w'}{\alpha}\binom{t'}{\beta}X^{\alpha+p\beta}Y^{(w'-\alpha)+p(t'-\beta)}.
\]
\unsure{Toujours un|Trouver un|Toujours est|Terme en} $X^iY^j$ \unsure{correspond|correspondant} à $\alpha=u,\beta=r$, il faudrait $u=\alpha\leq w'=u+v-p$ donc $v\geq p$, ce qui n'est pas, donc il n'y a pas de
% archive sheet number: 19

\scanpage{21}
tels termes. Donc $c_{i,j}\equiv0\ (p)$ dans ce cas, ce qui prouve (remplaçant $i,j$ par $u,v$) que $c_{u,v}\equiv0\ (p)$ donc encore $c_{ij}\equiv c_{u,v}\,c_{r,s}\ (p)$.

Le lemme 1 résulte du corollaire par une récurrence immédiate sur la longueur $e$ des développements $p$-adiques écrits de $u$ et de $v$.

\uline{Cor. 2} Soit $N$ un entier \inserted{$\geq0$} tel que $N+1\equiv0\ (p)$, posons $N+1=(N'+1)p$ donc $N=(N'+1)p-1=N'p+(p-1)$ \deleted{\unsure{$\in$|$($}} de sorte que à décalage de 1 près\unsure{\}} , \deleted{les} $N+1$ et $N'+1$ ont $=$ coeff. du développem\textsuperscript{t} $p$-adique, ainsi que $N$ et $N'$ (quand on fait abstraction du premier coeff. $p-1$ de $N$). Soit alors $\varepsilon\geq1$, on a
\[
\begin{aligned}
c_{N+1,p^\varepsilon-1}\deleted{\unsure{\not\equiv0\ (p)\iff}}&\equiv c_{N'+1,p^{\varepsilon-1}-1}\deleted{\unsure{\not\equiv0}}\ (p)\deleted{\unsure{\iff0}}\left\{\begin{aligned}&\equiv1\ (p)\ \text{si }N'+1\equiv0\ (p^{\varepsilon-1})\\&\equiv0\ (p)\ \text{si }N'+1\not\equiv0\ (p^{\varepsilon-1})\end{aligned}\right.\\
c_{N,p^\varepsilon}\deleted{\unsure{\not\equiv0\ (p)\iff}}&\equiv c_{N',p^{\varepsilon-1}}\deleted{\unsure{\not\equiv0}}\ (p)\deleted{\unsure{\iff}}\left\{\begin{aligned}&\equiv v_{\varepsilon-1}+1\ \inserted{(p)}\ \text{où}\\&\deleted{\unsure{v_\alpha}}\quad N'=\sum v_\alpha p^\alpha\end{aligned}\right)
\end{aligned}
\]

\uline{Dém.} Le développement $p$-adique de $p^\varepsilon-1$ est
\[
p^\varepsilon-1=(p-1)+(p-1)p+\cdots+(p-1)p^{\varepsilon-1}
\]
de m $p^{\varepsilon-1}-1=(p-1)+\cdots+(p-1)p^{\varepsilon-2}$ si $\varepsilon\geq2$ ($=0$ si $\varepsilon=1$).

Posant
\[
\left.\begin{aligned}
N'+1&=u_0+u_1p+\cdots+u_\nu p^\nu\\
N'&=v_0+v_1p+\cdots+v_\nu p^\nu
\end{aligned}\right\}\quad(\nu\geq\varepsilon-1)
\]
d'où
\[
\begin{aligned}
N+1&=0+u_0p+u_1p^{\unsure{2}}+\cdots+u_\nu p^{\nu+1}\\
N&=p-1+v_0p+v_1p^{\unsure{2}}+\cdots+v_\nu p^{\nu+1}.
\end{aligned}
\]
On trouve que
\[
\left\{\begin{aligned}
c_{N+1,p^\varepsilon-1}&\equiv c_{u_0,p-1}\,c_{u_1,p-1}\ \cdots\ c_{u_{\varepsilon-2},p-1}\quad(p)\qquad(\equiv1\ \inserted{(p)}\ \text{si }\varepsilon=1)\\
c_{N'+1,p^{\varepsilon-1}-1}&\equiv c_{u_0,p-1}\,c_{u_1,p-1}\ \cdots\ c_{u_{\varepsilon-2},p-1}\quad(p)\qquad(\equiv1\ (p)\ \text{si }\varepsilon=1)
\end{aligned}\right.
\]
\[
\left\{\begin{aligned}
c_{N,p^\varepsilon}&\equiv c_{v_{\varepsilon-1},1}\equiv v_{\varepsilon-1}+1\\
c_{N',p^{\varepsilon-1}}&\equiv c_{v_{\varepsilon-1},1}\equiv v_{\varepsilon-1}+1
\end{aligned}\right.
\]

\deleted{\unsure{Prop}}\uline{Lemme 2} Soit $\mathbb{Z}_p$ le localisé de $\mathbb{Z}$ en $p$, $E$ un $\mathbb{Z}_p$-module, $N\in\mathbb{N}$, $e\geq0$ tel que $p^e\geq N+1$ \deleted{\unsure{(donc $e\geq0$)}}, $I$ l'ens des entiers de la forme $N$ et $N+p^\varepsilon$, $0\leq\varepsilon\leq e$, $(x_i)_{i\in I}$ une famille d'éléments de $E$, tels qu'on ait
\[
\binom{j}{i}(x_j-x_i)=0\qquad i=N,N+1,\quad i<j\in I\quad(\text{donc }j=N+p^\varepsilon,\ 0\leq\varepsilon\leq e)
\]
\margin{\uline{Rmq} \unsure{même} si $E$ un $\mathbb{F}_p$-\unsure{es vectoriel|module}, on \unsure{ne} peut conclure que $x_i=x_N$ $\forall i\in I$ (ex : \inserted{$p=2$,} $N=2,e=2$, donc on a $x_2,x_3,x_4,x_6$ on a $x_2=x_3=x_6$ mais \uline{pas} néc. $=x_4$}% margin note written sideways
% archive sheet number: 20

\scanpage{22}
\deleted{\unsure{Lemme 3}} Alors on a $\ x_N=x_{N+1}$

\uline{Dém} Posons
\[
\left\{\begin{aligned}
&\deleted{\unsure{y_0=x_{N+1}-x_N}}\\
&y_\varepsilon=x_{N+p^\varepsilon}-x_{N\deleted{\unsure{+1}}}\qquad 0\leq\varepsilon\leq e.
\end{aligned}\right.
\]
d'où
\[
\begin{aligned}
&y_\varepsilon-y_0=x_{N+p^\varepsilon}-x_{N+1}\\
&\deleted{\unsure{y_\varepsilon=x_{N+p^\varepsilon}-x_{N+1}}}
\end{aligned}
\]
L'hypothèse s'écrit donc, en distinguant les cas $i=N,i=N+1$
\[
\begin{aligned}
&\deleted{\unsure{\binom{N+1}{N}y_0=0}}\hspace{5em}\deleted{\unsure{(i=N,j=N+1)}}\\
&(\deleted{\unsure{\bar S}}S)=S(N,e,(y_\varepsilon)_{0\leq\varepsilon\leq e})\left\{\begin{aligned}
&c_{N,p^\varepsilon}y_\varepsilon=0\qquad 0\leq\varepsilon\leq e\\
&c_{N+1,p^\varepsilon-1}(y_\varepsilon-y_0)=0\qquad 0\leq\varepsilon\leq e\quad\text{(le cas $\varepsilon=0$ étant trivial)}
\end{aligned}\right.
\end{aligned}
\]
La conclusion s'écrit $y_0=0$. En fait, nous allons déduire cette conclusion du système plus faible \inserted{(équivalent si $E$ est \deleted{un} annulé par $p$, i.e. est un $\mathbb{F}_p$-vectoriel)}
\[
(\bar S^{\unsure{\prime|{}}})=\bar S(N,e,(y_\varepsilon)_{0\leq\varepsilon\leq e}\left\{\begin{aligned}
&y_\varepsilon=0\quad\text{si }c_{N,p^\varepsilon}\not\equiv0\ (p)\qquad 0\leq\varepsilon\leq e\\
&y_\varepsilon-y_0=0\quad\text{si }c_{N+1,p^\varepsilon-1}\not\equiv0\ (p)\qquad 1\leq\varepsilon\leq e.
\end{aligned}\right.
\]

Pour prouver que ce système implique $y_0=0$, on récurre sur $N$. Comme la première équation pour $\varepsilon=0$ donne
\[
y_0=0\quad\text{si}\quad N+1\not\equiv0\ (p)
\]
la conclusion est triviale pour $N=0$. \deleted{\unsure{Alors on}} On suppose donc $N>0$, \inserted{et prouvé pour les $N'<N$.} La conclusion est triviale si $N+1\not\equiv0\ (p)$. Donc OPS
\[
N+1\equiv0\ (p)\qquad\text{soit}\qquad N+1=(N'+1)p
\]
où $N'=\frac{N+1}{p}-1<N$ (car $\frac{N+1}{p}<N+1$). Alors le cor. précédent montre que le syst. $\bar S=\bar S(N,e,(y_\varepsilon)_{0\leq\varepsilon\leq e})$ \inserted{(compte tenu de $N+1\equiv0\ (p)$ qui nous permet d'oublier la première équation pour $\varepsilon=0$)} équivaut au système (où $e'=e-1$, $y'_\varepsilon=y_{\varepsilon+1}$ pour $0\leq\varepsilon\leq e'=e-1$
\[
\deleted{\unsure{\bar S'=\bar S'(N',e',(y'_\varepsilon)_{1\leq\varepsilon}}}\quad
\left\{\begin{aligned}
&y'_\varepsilon=0\quad\text{si }c_{N',p^\varepsilon}\not\equiv0\ (p)\qquad 0\leq\varepsilon\unsure{=|\leq}e'\\
&y'_\varepsilon-y_0=0\quad\text{si }c_{N'+1,p^\varepsilon-1}\not\equiv0\ (p)\qquad 0\leq\varepsilon\leq e'
\end{aligned}\right.
\]
Pour $\varepsilon=0$ cette dernière équation donne \deleted{\unsure{$x$}}
\[
y'_0=y_0
\]
donc on peut se borner à $1\leq\varepsilon\leq e'$ (et se rappeler $y'_0=y_0$).
% archive sheet number: 21

\scanpage{23}
Mais alors ce système n'est autre que $\bar S'=\bar S'(N',e',(y'_\varepsilon)_{0\leq\varepsilon\leq e'})$. Notons que $p^e\geq N+1=(N'+1)p$ implique $p^{e-1}\geq N'+1$ i.e. $p^{e'}\geq N'+1$. Donc on est dans les $\overline{\mathrm{mes}}$ conditions que précédem\textsuperscript{t}, avec cette fois-ci $N'<N$. Par hyp. de récurrence on a donc $y'_0=0$, donc $y_0=0$, cqfd. Ouf !

\uline{Théorème 1} Soient $E$ un groupe abélien, $N\in\mathbb{N}$, $(x_i)_{i\in\mathbb{N}}$ une famille d'\unsure{éléments|élément} de $E$, telle qu'on ait
\[
\binom{j}{i}(x_j-x_i)=0\qquad\text{pour}\quad j>i\geq N
\]
Alors on a $\quad x_i=x_N\qquad\forall i\geq N$.

\uline{Dém} Quitte à localiser en tous les nbs premiers $p$, OPS $E$ un module sur \inserted{un} $\mathbb{Z}_p$. Il suffit de prouver $x_{N+1}=x_N$ (car par récurrence on \unsure{déduira|aura} de \unsure{m} $x_{N+1}=x_{N+2}$ etc). Choisissons $e\inserted{\in\mathbb{N}}$ tel que $p^e\geq N+1$, ne retenons que les $x_i$ pour $i$ de la forme $N$ ou $N+p^\varepsilon$, $0\leq\varepsilon\leq e$. Le lemme 2 nous montre que $x_{N+1}=x_N$, cqfd.

\uline{Cor 1} \unsure{Soient|Soit} $k$ un anneau commutatif, $S=k\{T\}$ l'alg. des polynômes à puiss. div., \deleted{$M$ un $k$-module\unsure{, $S$, $S^{+(N)}$}} \deleted{\unsure{$M$}} $S^{+(N)}$ l'idéal à p. div. \deleted{engendré par $T$ \unsure{soit|$i\geq N$}} défini par
\[
S^{+(N)}=\sum_{i\geq N}k.T^{(i)}.
\]
\deleted{et} Considérons, pour \unsure{tt|le|un} \deleted{\unsure{$M$}} $k$-module $M$, le $S$-module gradué
\[
\deleted{\unsure{=}}\ \unsure{w|\omega}_N(M)=M\otimes_S S^{+(N)}=\sum_{i\geq N}M.T^{(i)}
\]
Alors :
% archive sheet number: 22

\scanpage{24}
a) Le foncteur $\unsure{w|\omega}_N$ (de la catégorie des $k$-modules vers celle des $S$-modules gradués) est pleinement fidèle.

b) \inserted{$\forall M$,} l'application
\[
\deleted{\unsure{x}}M\xrightarrow{\varphi}\prod_{i\geq0}w_N(M)_i=\prod_{i\geq N}M.T^{(i)}
\]
donnée par
\[
\varphi(x)=(x.T^{(i)})_{i\geq N}
\]
induit une bijection de $M$ avec le sous-ens. de $\prod_{i\geq N}\deleted{\unsure{M}}w_N(M)_i$ formé des systèmes $(x_i)_{i\geq N}$ satisfaisant la condition
\[
T^{(j)}x_i=c_{ij}x_{i+j}\qquad\forall i\geq N,\ j\geq1\tag*{$(*)$}
\]

\uline{Dém}

\deleted{a)}b) $x_i=xT^{(i)}$, on a $T^{(j)}x_i=xT^{(i)}T^{(j)}=c_{ij}xT^{(i+j)}=c_{ij}x_{i+j}$, donc $\varphi(x)$ se trouve dans le sous-ens indiqué. L'injectivité de $\varphi$ est claire. Bijectivité ? Posons $x_i=\xi_i T^{(i)}$, $\xi_i\in M$. La condition $(*)$ s'écrit
\[
\begin{aligned}
\xi_i T^{(i)}T^{(j)}&=c_{ij}\xi_{i+j}T^{(i+j)}\\[-2pt]
\shortparallel&\\[-2pt]
c_{ij}\xi_iT^{(i+j)}&
\end{aligned}
\]
i.e.
\[
c_{ij}(\xi_{i+j}-\xi_i)=0\qquad\text{pour}\quad i\geq N,\ j\geq1,
\]
et il résulte du th. que les $\xi_i$ sont tous égaux à un $=$ $\xi\in M$, cqfd.

a) \unsure{Soient|Soit} $M,M'$ sur $k$. Alors $\operatorname{Hom}(w_N(M),w_N(M'))$ s'identifie à l'ens des systèmes de $k$-hom $(i\geq N)$
\diagram{\begin{tikzcd}[ampersand replacement=\&,column sep=large,row sep=small]
u_i\colon M \arrow[r,phantom,"\unsure{-|\to}"] \arrow[d,phantom,"\rotatebox{90}{$\simeq$}"] \& M' \arrow[d,phantom,"\rotatebox{90}{$\simeq$}"]\\
M.T^{(i)} \arrow[d,equal] \& M'.T^{(i)} \arrow[d,equal]\\
w_N(M)_i \& w_N(M')_i
\end{tikzcd}}
% archive sheet number: 23

\scanpage{25}
tels que $\forall j\geq1$, on ait comm. dans
\diagram{\begin{tikzcd}[ampersand replacement=\&,column sep=large,row sep=large]
M.T^{(i)} \arrow[r,"{.T^{(j)}}"] \arrow[d,"{u_i\otimes T^{(\unsure{i|h})}}"'] \& M T^{(i+j)} \arrow[d,"{u_{i+j}\otimes T^{(\unsure{i|h}+j)}}"]\\
M'.T^{(i)} \arrow[r,"{.T^{(\unsure{j|i})}}"] \& M'.T^{(i+j)}
\end{tikzcd}}
i.e. tels qu'on ait
\[
c_{ij}(u_{i+j}-u_i)=0\qquad\text{si}\quad i\geq N,\ j\geq1,
\]
et le th. nous permet de conclure que tous les $u_i$ sont égaux, d'où la conclusion
\[
\operatorname{Hom}_k(M,M')\xrightarrow{\sim}\operatorname{Hom}_{S\text{-mod gr.}\unsure{d|s|{}}}(\unsure{w|\omega}_N(M),w_N(M'))
\]

\uline{Rmq 1} Le foncteur $w_N$ admet un adjoint à droite $\pi_N$
\[
\operatorname{Hom}_{S\text{-mod gr.}d}(w_N(M),\deleted{\unsure{M'}}\mathcal M)\unsure{\simeq|\xrightarrow{\sim}}\operatorname{Hom}_k(M,\pi_N(\mathcal M))
\]
en posant
\[
\begin{aligned}
\pi_N(\mathcal M)&\subset\prod_{i\geq N}\mathcal M_i\\
\shortparallel&\\
\bigl\{(x_i)_{i\geq N}\ &\big|\ x_iT^{(j)}=c_{ij}x_{j+j}\quad\forall i\geq N,\ j\geq1\bigr\}
\end{aligned}
\]
En considération de ceci, \deleted{il existe} les parties a) et b) du corollaire 1 sont équivalentes, b) signifiant que l'hom can. d'adjonction
\[
M\longrightarrow\pi_N(w_N(M))
\]
est \deleted{\unsure{bij}} un isomorphisme.

\uline{Rmq 2} Considérons $w_0(M)=M\otimes_k S=M\otimes_k k\{T\}$, alors $w_N(M)=S^{+(N)}.w_0(M)$. \deleted{Donc \unsure{pour} en un sens, $w_0$ est} \inserted{On peut aussi écrire, si $N\geq N'\geq0$}
\[
w_N(M)\simeq\deleted{\unsure{\prod}}\sum_{i\geq N}w_{N'}(M)_i\overset{\mathrm{def}}{=}\underset{\text{(``troncature'')}}{\tau_{\geq N}}(w_{N'}(M)).
\]
Donc en
% archive sheet number: 24

\scanpage{26}
un sens les foncteurs ``ombres'' $\unsure{w|\omega}_N$ sont d'autant plus grossiers que $N$ est plus grand. Néanmoins ils sont tous pleinement fidèles (\uline{NB} les foncteurs $\tau_N$ ne sont pas pl. fidèles bien sur, pas $=$ sur la sous-catégorie des ``ombres virtuelles d'ordre $N-1$'', i.e. des modules gradués dont les composantes sont nulles en degrés $\leq N-1$.

Soit $\operatorname{Omb}(k)$ la catégorie des ``ombres virtuelles'' de $k$-modules, i.e des $\deleted{\unsure{C}}S=k\{T\}$-modules gradués [à degrés $\geq0$ ?]. Soit, $\forall N\geq0$, $\operatorname{Omb}_{N\text{-}\mathrm{tr}}(k)$ la sous-catégorie pleine des modules $N$-Négligeables, i.e. tels que $\mathcal M_i=0$ si $i\geq N$. C'est une sous-catégorie épaisse, la catégorie quotient est can. équivalente \deleted{à celle} \inserted{à la ss-cat. pleine de (Omb) formée} des ``\inserted{$N$-}Ombres virtuelles'', i.e. des $\deleted{\unsure{k}}S$-modules gradués tels que $\mathcal M_i=0$ si $i<N$.

Le foncteur passage au quotient s'identifie alors au foncteur ``troncature'' $\tau_N$. On notera que
\[
\tau_N(\mathcal M)\lhook\joinrel\longrightarrow\mathcal M
\]
(inclusion fonctorielle) \inserted{correspondant au fait que} le foncteur \deleted{\unsure{quo}} passage au quotient admet un adjoint à droite (p. fid. comme de juste)
\[
\operatorname{Hom}_{\mathrm{Omb}}(\mathcal N,\tau_N(\mathcal M))\simeq\operatorname{Hom}_{\substack{N\text{-}\mathrm{Omb}\\\text{ou }\mathrm{Omb}}}(\tau_N(\unsure{\mathcal M|\mathcal N}),\tau_N(\mathcal M)).
\]

Appelons ``négligeable'' \unsure{une ombre|toute ombre|les ombres} $\mathcal M$ qui est $\underset{\longrightarrow}{\lim}$ \inserted{filtrée} d'ombres $N$-négligeables (pour $N$ variable), i.e. \unsure{telle|telles} que \inserted{$\forall i\in\mathbb{Z}$,} $\forall x\in\mathcal M_i$, $\exists N\in\mathbb{N}$ tel que $T^{(\alpha)}.x=0$ pour $\alpha\geq N$ i.e. $S^{+(N)}.x=0$. C'est encore une sous-catégorie épaisse, et le localisé est
% archive sheet number: 25

\scanpage{27}
$S=k\{T\}$ (polynômes à p.d.) anneau gradué à degrés $\geq0$

$M$ un $k$-module, $N\geq0$, d'où $M\otimes_k S$ et par suite
\[
M\otimes_k S\longrightarrow\varphi_N\tau_N(M\otimes_k S)
\]
\uline{Thm} L'hom précédent est un isom en degrés $\geq0$, i.e. on a
\[
M\otimes_k S\xrightarrow{\sim}\varphi_{0N}\tau_N(M\otimes_k S)\tag*{(\unsure{$1$|$\alpha$|$*$})}
\]
Enoncé équivalent :
\[
M\otimes_k S\in\operatorname{Im\,ess}\varphi_{0N}\qquad(\forall N\geq0)
\]
ou encore
\[
M\otimes_k S\in\operatorname{Im\,ess}\varphi_{0,\unsure{\infty|\omega}}
\]

\uline{NB} En degrés $\geq N$, c'est clair que \unsure{$\alpha$|ce|c'} est un isom, il suffit donc de regarder les degrés $0\leq i<N$. Le cas $i=0$ a été traité déjà par une démonstration spéciale (on \unsure{la|le} retrouvera dans la dém. qui va suivre).

Reformulation du \uline{Thm}

Si $0\leq i\leq N$, on a
\diagram{\begin{tikzcd}[ampersand replacement=\&,column sep=small,row sep=small]
\left(\begin{smallmatrix}\xi T^{(i)}\\ \uparrow\\ \xi\end{smallmatrix}\right) \& (M\otimes_k\deleted{\unsure{B}}S)_i \arrow[r,"\sim"] \arrow[d,phantom,"\rotatebox{90}{$\simeq$}"] \& \operatorname{Hom}^{i}\bigl(\tau_{N-i}(S),M\otimes_k S\bigr)\\
\& M \&
\end{tikzcd}}
(seul le cas $0\leq i\unsure{<|\leq}N$ doit être établi)

\uline{Cor} $\operatorname{Ext}^n\bigl(\tau_{N-i}(S)[-i],M\otimes_k S\bigr)=0$ si $n>0$ $\forall i\geq0$ (seul le cas $i<N$ est non trivial\deleted{\unsure{ici}})

Ce théorème équivaut à l'énoncé suivant, où $k$ ne figure plus

\uline{Th} Soient $0\leq i\leq N$ deux entiers, $M$ un $\mathbb{Z}$-module, $(\xi_k)_{k\geq N}$ une famille d'éléments de $M$, telle que l'on ait
\[
\unsure{e|c}_{l,k}\xi_k=c_{l,k-i}\xi_{k+l}\qquad\text{pour }k\geq N,\ l\geq0\tag*{(1)}
\]
% archive sheet number: 26

\scanpage{28}
Alors il existe un unique $\xi\in M$ tel que l'on ait
\[
\xi_k=c_{\deleted{\unsure{k}}i,k-i}\xi\qquad\left(=\binom{k}{i}\xi\right)\qquad\text{pour tt }k\geq N\tag*{(2)}
\]
(\uline{NB} $c_{k,l}$ désigne le coefficient binômial $\frac{(k+l)!}{k!\,l!}$)

\uline{Dém} \textcircled{1} Réduction au cas où $\exists p$ nb premier, tel que $\deleted{\unsure{p}}n\,\mathrm{id}_M$ bijectif $\forall n\in\mathbb Z$ tel que $(n,p)=1$, i.e. $M$ un $\mathbb Z_p$ module ($\mathbb Z_p\overset{\mathrm{def}}{=}\mathbb Z_{(p\mathbb Z)}$).

Supposons que ce soit prouvé dans ce cas, et qu'on sache qu'il existe alors $k=k(p)$ tel que $\deleted{\unsure{\xi}}c_{\deleted{\unsure{k}}i,k-i}\not\equiv0\ (p)$ (i.e $c_{k,k-i}\in\mathbb Z_p^{\unsure{\times|*}}$]. Prouvons le th. en général. Les $c_{k,k-i}$ ($k\geq N$) ne sont pas divisibles par un $=$ nb premier, donc sont premiers dans leur ensemble, donc $\exists k_1\ldots k_r\geq N$ et des entiers $m_1,\ldots m_r\in\mathbb Z$ tels que
\[
\sum_{\alpha=1}^{r}m_\alpha c_{k_\alpha,k_\alpha-i}=\unsure{1|d}.
\]
Donc les relations (2) impliquent, en multipliant par les $m_\alpha$ et ajoutant
\[
\xi=\sum_1^r m_\alpha\xi_{k_\alpha}
\]
d'où l'unicité de $\xi$. Prouvons qu'avec ce choix de $\xi$, on a bien (2) i.e.
\[
\xi_k=c_{k,k-i}\sum_1^{\unsure{r}} m_\alpha\xi_{k_\alpha}
\]
Pour prouver que cette relation est vraie, \deleted{\unsure{p|k ?}} il suffit de prouver qu'elle est vraie dans chaque $M_p=M\otimes_{\mathbb Z}\mathbb Z_p$ ($p\in\mathbb P$). Mais c'est une conséquence immédiate du fait que le th. est vrai pour $M_p$.
% archive sheet number: 27

\scanpage{29}
\textcircled{2} Cas où $M$ est un $\mathbb Z_p$-module

Notons d'abord qu'il existe des $k$ arbitrairem\textsuperscript{t} grands tels que $c_{i,k-i}\in\mathbb Z_p^{\unsure{\times|*}}$ i.e. $c_{i,k-i}\not\equiv0\ (p)$. Posant $k-i=j$, il suffit de prendre $j=p^\varepsilon$, avec $\varepsilon$ tel que $p^\varepsilon>i$, donc $k=i+p^\varepsilon$. Cela montre l'unicité de $\xi$, car si $k_0$ est tel que $c_{i,k_0-i}\in\mathbb Z_p^\times$, on aura
\[
\deleted{\unsure{(2)}}\qquad \xi=c_{i,k_0-i}^{-1}\xi_{k_0}
\]
Considérons alors le système des
\[
\deleted{\unsure{(3)}}\qquad \xi'_k=c_{i,k-i}\xi.\quad,
\]
on voit de suite que \deleted{\unsure{ce}} $(\xi'_k)_{k\geq N}$ satisfait \deleted{\unsure{aux hyp|analogue}} les conditions (1)\unsure{,|{}} (où on substitue $\xi'_k$ à $\xi_k$), et $\xi'_{k_0}=\xi_{k_0}$, donc posons
\[
\unsure{\eta|y}_k=\xi'_k-\xi_k,
\]
\deleted{les} $(\eta_k)_{k\geq N}$ satisfait encore (1), et de plus $\eta_{k_0}=0$. Il suffit de prouver alors $\eta_k=0$, \deleted{donc} $\forall k\geq N$. Donc on est ramené au

\uline{Corollaire} Si $(\xi_k)_{k\geq N}$ \deleted{est} satisfait (1), et \inserted{si} de plus \deleted{\unsure{$(1\frac12)$}} on \unsure{donné|donne} un $k_0\geq N$ tel que
\[
\xi_{k_0}=0\tag*{$(1')$}
\]
\[
c_{i,k_0-i}\not\equiv0\ (p)\tag*{(3)}
\]
alors $\qquad \xi_k=0\qquad\forall k\geq N$.
% archive sheet number: 28

\scanpage{30}
\uline{Dém} En deux étapes

\uline{\unsure{Lemme|Lem}, 1}, $\xi_k=0$ si \inserted{$k\geq N$ tel que} $c_{i,k-i}\not\equiv0\ (p)$.

\uline{Lemme} 1.1. Soient $\deleted{\unsure{k',k|k_1,k}}\ k'\geq k\geq N$ tels que, posant $j=k-i$, $j'=k'-i\deleted{\unsure{\prime}}$, on ait
\[
c_{i,j}\not\equiv0\ (p)\quad\text{et}\quad c_{i,j'}\not\equiv0\ (p)\tag*{(4)}
\]
Alors on a l'équivalence
\[
c_{j'-j,j}\not\equiv0\ (p)\iff c_{j'-j,j+\deleted{\unsure{j}}i}\not\equiv0\ (p)\tag*{(5)}
\]
et si les conditions équivalentes (5) sont satisfaites, on a l'équivalence
\[
\xi_k=0\iff\xi_{k'}=0.\tag*{(6)}
\]

\uline{Dém de} 1.1. Calculons dans $\mathbb Q$
\[
\begin{aligned}
\text{(a)}\quad &c_{j'-j,j+i}\,/\,c_{j'-j,j}=\frac{(j'+i)!}{\deleted{(j'-j)!}(j+i)!}:\frac{(j'\deleted{\unsure{+j-j}})!}{\deleted{(j'-j)!}j!}\\
&\text{et }\deleted{\unsure{/}}\\
\text{(b)}\quad &c_{i,j'}\,/\,c_{ij}=\frac{(i+j')!}{\deleted{i!}j'!}:\frac{(i+j)!}{\deleted{i!}j!}
\end{aligned}
\]
\deleted{\unsure{et}}
\deleted{\[
\text{(b)}\quad c_{i,j'}/c_{ij}=\frac{(i+j')!}{i!\,j'!}:\frac{(i+j)!}{i!\,\unsure{j!}}
\]}
\deleted{Or dans $\mathbb Z_p$, pour la \unsure{rel.|val.} de divisibilité, on \unsure{le quotient de (a) et (b) est} grâce à (4),
\[
(i+j)!\sim i!\,j!\qquad (i+j')!\simeq i!\,j'!
\]
donc le premier \unsure{membre} de (a) est \unsure{eq} à
\[
\frac{i!\,j'!}{i!\,j!}:\frac{(j'+i-j)!}{i!}=
\]}
Donc les deux \unsure{quantités|quotients} (a) et (b) sont égaux, et comme (b) est un quotient de deux unités \deleted{$p$-adiques} \inserted{de $\mathbb Z_p$}, on voit que $c_{j'-j,j+i}$ et $c_{j'-j,j}$ sont équivalents pour la divisibilité, donc l'un est une unité sss l'autre l'est. Quand il en est ainsi, la
% archive sheet number: 29

\scanpage{31}
relation $(1)$ avec $l=k'-k=j'-j$ donne
\[
c_{j'-j,j+i}\,\xi_k=c_{j'-j,j}\,\xi_{k'\deleted{\unsure{+l|+i}}}
\]
et comme les coeffts sont des unités de $\mathbb Z_p$, on en conclut que $\xi_k=0\Longleftrightarrow\xi_{k'}=0$, cqfd.

\uline{Dém du lemme} 1\quad Considérons le développem$^{\rm t}$ $p$-adique de $i$
\[
 i=i_0+i_1p+\cdots+i_rp^r \qquad 0\leq i_\alpha\leq p-1\quad(0\leq\alpha\leq r)
\]
où on prendra $r$ arbitrairement grand (on ne suppose donc \uline{pas} que $i_r\neq0$ !). Considérons de $=$ le développem$^{\rm t}$ de $j=k-i$
\[
 j=j_0+j_1p+\cdots+j_rp^r\qquad 0\leq j_\alpha\leq p-1\quad(0\leq\alpha\leq r)
\]
L'hyp. $c_{i,j}\not\equiv0\ (p)$ signifie que
\[
 i_\alpha+j_\alpha\leq p-1\qquad\forall\ 0\leq\alpha\leq r
\]
Soit \deleted{\unsure{$\delta_\alpha$}}
\[
 \delta_\alpha=(p-1)-(i_\alpha+j_\alpha)\qquad 0\leq\delta_\alpha\leq p-1
\]
i.e.
\[
 j'_\alpha\overset{\text{déf}}{=}j_\alpha+\delta_\alpha=(p-1)-i_\alpha
\]
donc
\[
 0\leq j'_\alpha\leq p-1,
\]
et considérons
\[
 j'=\sum_0^r j'_\alpha p^\alpha
\]
donc $j'\geq j\ (\geq N)$ et
\[
 j'-j=\delta\overset{\text{déf}}{=}\sum_0^r\delta_\alpha p^\alpha.
\]
% margin note written sideways
\margin{\uline{Présentation plus simple} Soit $r\geq0$ tel que $p^r>\deleted{\unsure{\geq j|j}}i+j$ soit $\deleted{\unsure{\delta=p^r-1-}}\,\delta=p^r-1-(i+j)$ d'où $i+j+\delta=p^r-1$, $i+(j+\delta)=''$. Alors on a $\left\{\begin{aligned}c_{i,j+\delta}&\not\equiv0\ (p)\\c_{i+j,\delta}&\not\equiv0\ (p)\end{aligned}\right.$ (en plus de $c_{i,j}\not\equiv0\ (p)$) donc $\xi_k=0\Longleftrightarrow\xi_{k'}=0$, où $k'=j'+i=j+\delta+i=p^r-1$.}
On a alors \inserted{(en plus de $c_{i,j}\not\equiv0$)}
\[
\begin{aligned}
 c_{i,j'}&\not\equiv0\ (p) &&\text{car }i_\alpha+j'_\alpha\underset{\forall\alpha}{(=p-1)}\leq p-1\\
 c_{j'-j,j}&\deleted{\unsure{\not\equiv0|=}}=c_{\delta,j}\not\equiv0\ (p) &&\text{car }\delta_\alpha+j_\alpha=j'_\alpha\leq p-1
\end{aligned}
\]
% archive sheet number: 30

\scanpage{32}
donc en vertu de $(1.1)$ on a
\[
 \xi_k=0\Longleftrightarrow\xi_{k'}=0
\]
où
\[
 k'=i+j'=\sum_0^r(p-1)\,p^\alpha=p^{r+1}-1,
\]
donc on trouve que
\[
 \xi_k=0\Longleftrightarrow\xi_{p^{r+1}-1}=0.
\]
Pour la $=$ raison, on trouve, si $\widetilde k\deleted{\unsure{=i+j'|=i+j}}\inserted{=i+\widetilde j}$ satisfait la $=$ condition que $k$, savoir $c_{i,\unsure{k'|\widetilde k}-i}=0$
\[
 \xi_{\widetilde k}=0\Longleftrightarrow\xi_{p^{r+1}-1}=0
\]
où on choisit $r$ tel que $p^{r+1}\geq\deleted{\unsure{k|i}}j,\widetilde j\ (\geq i)$.

En comparant, on trouve
\[
 \xi_k=0\Longleftrightarrow\xi_{\widetilde k}=0
\]
i.e. si \uline{un} $\xi_k$ avec $k\geq N$, $c_{k,k-i}\not\equiv0\ (p)$ est nul, tous le sont, cqfd.

\uline{Fin de la démonstration} Il \deleted{suffit de}\inserted{faut} prouver $\xi_k=0\quad\forall k\geq N$. On a
\[
 c_{l,k}\,\xi_k=c_{l,k-i}\,\xi_{k+l}\qquad\forall l\geq0
\]
et il suffit de prouver qu'il existe $l\geq0$ tel que $c_{l,k}\not\equiv0\ (p)$, et de plus $\xi_{k+l}=0$, ce qui est le cas \deleted{si} en vertu de $1.1.$ si $c_{i,(k+l)-i}\not\equiv0\ (p)$. Posons encore $k-i=j$ i.e. $k=i+j$, on veut
% archive sheet number: 31

